\chapter*{A field guide to the next problem}
\markboth{A field guide to the next problem}{A field guide to the next problem}
\addcontentsline{toc}{chapter}{A field guide to the next problem}
The sum $1^3+\cdots+n^3$ equals $(1+\cdots+n)^2$ for every positive integer $n$. Chapter~\ref{ch:sums} proves that sentence. The chance that $n$ hats come back with nobody holding their own tends to $1/e$. Chapter~\ref{ch:probability} proves that sentence, from the derangement count and Taylor's formula for the exponential. One law is an identity. The other is a proportion. A list of checks is neither.

The next problem will not say which of the two it is, and it will not name a chapter. Write the sentence it is asking. Name the objects, the domain, and the words ``every'' and ``some''. Then try a number you can compute by hand. If that number breaks the sentence, stop: you have a counterexample, and the word that was too large is usually ``every''. If the small numbers agree, keep the agreement. It is a suspect, not a verdict. The sums chapter already has a polynomial that matches $2^n$ at five points and equals $31$, not $32$, at the sixth. Forty agreements can fail too. A proof has to say why the agreement cannot stop.

Two problems follow. The first is a formula that produces primes for a long time, and then cannot. The second is a count of words on a circle, where the useful language changes halfway through. The same questions apply to both, and to a problem you have not seen yet.

\section*{Forty primes, and then a square}
Define $f(n)=n^2+n+41$ for integers $n\ge0$. The first values are
\[
f(0)=41,\quad f(1)=43,\quad f(2)=47,\quad f(3)=53,\quad f(4)=61.
\]
Each of these is prime. For $61$, the primes up to $\sqrt{61}<8$ are $2,3,5,7$, and none of them divides $61$. The same check, with a shorter list of trial primes, handles $41,43,47,53$.

The streak is much longer than five. In a letter extract printed in the \emph{Nouveaux M\'emoires de l'Acad\'emie royale des sciences de Berlin} for $1772$, pages $35$--$36$ (the volume appeared in $1774$; Enestr\"om $461$; \emph{Opera Omnia}, series~$1$, volume~$3$, pages $335$--$337$), Euler points to the progression that begins $41,43,47,53$, with general term $41-x+x^2$. He says that its first forty terms are prime.

Those forty terms are the values of $f$. For each integer $x$,
\begin{align*}
x^2-x+41
&=(x^2-2x+1)+(x-1)+41\\
&=(x-1)^2+(x-1)+41.
\end{align*}
Set $n=x-1$. As $x$ runs from $1$ to $40$, the index $n$ runs from $0$ to $39$, and the general term is $f(n)$. Euler's fortieth term is $f(39)=1601$. The next term of his progression is the value at $x=41$, which is $f(40)$.

\begin{examplex}[the term after the forty]
$f(40)=40^2+40+41=1681$, and $41^2=1681$. The factor is visible before you multiply:
\[
f(40)=40\cdot41+41=41(40+1)=41^2.
\]
The same rearrangement at the next multiple gives
\[
f(41)=41\cdot42+41=41\cdot43=1763.
\]
Both numbers are divisible by $41$ and larger than $41$, so neither is prime. Figure~\ref{fig:field-euler} marks the break at $n=40$.
\end{examplex}

\begin{figure}[ht]
\centering
\fig{field-euler}
\caption{Values of $f(n)=n^2+n+41$. The five small charcoal values are prime by the trial check in the text, and $n=39$ is the index of Euler's fortieth term. The rust value is the next term, $f(40)=41^2$.}
\label{fig:field-euler}
\end{figure}

Forty recorded primes are a long list. They are still a list. The identity $f(40)=41^2$ does not need those forty values: it is arithmetic. One further family comes from the single value $f(0)=41$. For every integer $k\ge1$,
\[
f(41k)=(41k)^2+41k+41=41(41k^2+k+1),
\]
and the second factor is an integer at least $41+1+1=43$. So $f(41k)$ is composite for every $k\ge1$. This is an infinite family of composite values. The square at $n=40$ is a smaller failure, found by computing the next term. The family $n=41k$ is a different list.

\begin{think}
Which integers $n$ make $41$ divide $f(n)$? Reduce $f(n)=n(n+1)+41$ modulo $41$. The two residue classes you find contain $n=0$, where $f(0)=41$, the failures $n=40$ and $n=41$, and every later integer in those classes.
\end{think}

Modulo $41$, the constant $41$ disappears, so $f(n)\equiv n(n+1)\pmod{41}$. The product $n(n+1)$ is $0$ modulo $41$ when $n\equiv0$ or $n\equiv-1$. Thus $41$ divides $f(n)$ for every $n$ in those two classes. Whenever that value is also larger than $41$, it is composite. This is why the prime value $f(0)=41$ forces later terms, on those two progressions, that are not prime. The next argument drops the number $41$ and keeps the pattern: a non-constant polynomial cannot stay prime for every $n$.

\section*{Every such formula breaks}
A constant polynomial can take prime values forever: the constant $41$ is prime at every $n$. The moment the polynomial is allowed to grow, the streak of primes has to break, and it has to break infinitely often. The statement uses the definition from Chapter~\ref{ch:sums}: an integer $p>1$ is prime when its only positive divisors are $1$ and $p$.

\begin{theoremx}[prime values of a polynomial]
Let $f$ be a polynomial with integer coefficients, of degree at least $1$. There are infinitely many integers $n\ge0$ for which $f(n)$ is not prime.
\end{theoremx}

\begin{proof}
The proof has three stages. First, differences of integer polynomials stay divisible by the gap in the input. Second, one value with absolute value at least $2$ produces a whole arithmetic progression of multiples of that absolute value. Third, a non-constant polynomial can hit any fixed integer only finitely often, so only finitely many terms of that progression can be prime.

\emph{The difference.} Let $f(X)=\sum_{j=0}^{d}c_jX^j$ with each $c_j$ an integer and $d\ge1$, so $c_d\ne0$. For integers $s,t$ and an integer $j\ge1$,
\[
s^j-t^j=(s-t)\sum_{i=0}^{j-1}s^{j-1-i}t^i.
\]
The sum is an integer, so $s-t$ divides $s^j-t^j$. The case $j=0$ is the difference $0$, which $s-t$ also divides. Therefore $s-t$ divides
\[
f(s)-f(t)=\sum_{j=0}^{d}c_j(s^j-t^j).
\]

\emph{An arithmetic progression of multiples.} If $|f(n)|\le1$ for every integer $n\ge0$, then each value lies in $\{-1,0,1\}$, and none of those three integers is prime. In that case every $n\ge0$ is already a failure, and there is nothing left to prove. Suppose instead that some integer $a\ge0$ has $m=|f(a)|\ge2$. Then $m$ divides $f(a)$. For each integer $k\ge0$ set $n_k=a+km$. The gap $n_k-a=km$ divides $f(n_k)-f(a)$, and $m$ divides the gap, so $m$ divides the difference. Hence $m$ divides $f(n_k)$. Also $n_k\ge0$.

\emph{Only finitely many of them can be prime.} Suppose $f(n_k)$ is prime. By the definition, that value is an integer $p>1$, and the only positive divisor of $p$ that is at least $2$ is $p$ itself. Here $m\ge2$ divides $f(n_k)$, so $f(n_k)=m$, and in particular $m$ itself is prime. Consider the polynomial $f(X)-m$ in $\Q[X]$. It is not the zero polynomial: if it were, $f$ would be the constant $m$, whose degree is $0$, not at least $1$. Proposition~\ref{prop:rootsfactor} says that a non-zero polynomial of degree $d$ has at most $d$ roots in $\Q$. The numbers $n_k=a+km$ for $k=0,1,2,\dots$ are distinct. Only finitely many of them can be roots of $f(X)-m$.

Thus $f(n_k)=m$ for only finitely many integers $k\ge0$. For every other such $k$, the integer $f(n_k)$ is divisible by $m\ge2$ and is not equal to $m$, so it is not prime. There are infinitely many such $k$.
\end{proof}

Check the theorem on $f(n)=n^2+n+41$ with $a=0$ and $m=41$. The progression is $n_k=41k$, and the factored form above is the general argument specialised to this polynomial. Check a constant, which the degree hypothesis excludes: $f(n)=41$ is prime for every $n$, and the proof never starts, because a constant of degree $0$ is not allowed into the statement. Check a polynomial that is never prime at all, such as $f(n)=-n-2$. Every value is negative, hence not prime, and the opening case $|f(n)|\le1$ does not apply. The progression built from $a=0$, where $f(0)=-2$ and $m=2$, still supplies infinitely many non-primes. The degree was the hypothesis that forced a prime streak to end.

No choice of integer coefficients, of any degree, gives a polynomial that is prime at every integer $n\ge0$, unless the polynomial is a constant prime. Euler's forty values are a long finite check. Changing the coefficients does not escape the theorem. That is the theorem: no non-constant polynomial with integer coefficients is prime at every integer $n\ge0$. This volume stops before the theorems that describe the primes. The words on a circle are the same standard, used on a count.

\section*{Questions to ask first}
The sentence ``$f(n)$ is prime for every $n\ge0$'' was larger than the forty values. The failure at $n=40$ was easy to see once it was computed. Many problems do not announce the breaking value. Before you reach for a new formula, name the object the problem gave you, and the chapter that already handles an object of that shape. Each question below is one such test: a chapter of this volume, and a small thing to compute before the general argument.

\begin{description}
\item[The sentence.] Write the domain and the quantifiers before any algebra. Over the reals, $x^2=1$ has two solutions; over the positive reals it has one. Chapter~\ref{ch:language} is that distinction. If you cannot say what would count as a counterexample, the sentence is not finished.
\item[A small case.] Compute the first value, the empty object, and one value where the formula might lose an endpoint. The binomial identity $\binom n0=1$, and the empty subset counted by $f_0=1$ in the circular-word argument below, are the same kind of check. A formula that fails at $n=0$ or at $n=2$ is the wrong count, whatever it does at $n=100$.
\item[A sum.] Write the first term and the last term. Look for a difference $F(k+1)-F(k)$, a second way to count the same set, or the binomial basis of Chapter~\ref{ch:sums}. When you shift an index, write the new range beside the old one. The term that falls off the end is the usual mistake.
\item[A count.] Say whether the objects are ordered, and whether a symmetry has already identified some of them. Chapter~\ref{ch:counting} removes a compulsory gap by a bijection, and it counts the complement when the forbidden configurations overlap. Add the empty configuration on purpose.
\item[Every positive integer.] If the claim is ``for every $n$'', decide what the inductive step is allowed to use. A step that needs two previous cases needs two base cases. If the obvious statement will not push forward, prove a stronger one that carries the extra fact. In Chapter~\ref{ch:sums} the bound $\sum_{k=1}^n 1/k^2<2$ does not survive the inductive step, while $\sum_{k=1}^n 1/k^2\le 2-1/n$ does.
\item[A sequence or a series.] Prove that the terms stay where the expression is defined. Then look for monotonicity, a closed interval the terms cannot leave, or a Cauchy estimate. Chapter~\ref{ch:sequences} records the trap in both directions: the terms of a convergent series tend to $0$, and a series whose terms tend to $0$ need not converge. If $f$ is continuous, solving $L=f(L)$ names the only possible limits of a recursion $u_{n+1}=f(u_n)$. It does not prove that the sequence converges.
\item[A comparison of sizes.] Track the sign of every denominator. If both sides scale by the same power, you may set that scale to $1$. Find the equality case before you expand anything: equal variables, or a vector parallel to another, often tell you which square to complete. Chapter~\ref{ch:inequalities} is that search.
\item[A polynomial identity.] Bound the degree of the difference, then exhibit more roots than that degree allows, unless the difference is the zero polynomial. That root count is Proposition~\ref{prop:rootsfactor}, and the coefficient ring has to be a field. Over $\Z/8\Z$ the polynomial $X^2-1$ has four roots, namely $1,3,5,7$, because that ring has zero divisors and the proposition does not apply. The prime-value argument used the proposition only in $\Q[X]$, where it does apply, and it used divisibility of integer values for the factor $m$.
\item[A complex number.] Decide whether you are adding or multiplying. Addition wants coordinates. Multiplication wants a modulus and an argument, and the argument is defined only up to multiples of $2\pi$. A root has as many values as its order, one for each argument in a full turn. Chapter~\ref{ch:complex} keeps those choices separate.
\item[An integral or a Taylor approximation.] Say whether the interval is closed and bounded, or whether an endpoint is infinite or blows up. Make the substitution on a truncated interval, and only then take the limit. A Taylor polynomial with a named remainder, from Chapter~\ref{ch:calculus}, often replaces an interchange you do not yet have a theorem for.
\item[A matrix.] Ask for the kernel, the image, and any subspace the matrix carries to itself, before you expand a determinant. A change of basis changes the columns of numbers and does not change the linear map. An orthogonal projection is the nearest point in a subspace; the residual vector is the certificate. Theorem~\ref{thm:spectral} says that a real symmetric matrix has an orthonormal basis of eigenvectors.
\item[A probability.] Name the outcomes, the probability of each, and what information was revealed. Chapter~\ref{ch:probability} is strict about the order of those steps. Add expectations whenever the sum is finite. Multiply probabilities only when the factors are independent, or when you have written a conditional probability on purpose. The hat check is a count of permutations. It is not a product of $n$ guesses.
\item[An order.] If the problem says that something must line up, look for a chain or an antichain, or for a quantity that drops at every legal step. In a well-founded order, as defined in Chapter~\ref{ch:orders}, a strictly decreasing sequence cannot be infinite. That is the test that a process stops.
\end{description}

None of these replaces a proof. They tell you which proof to attempt, and which tempting step is not yet justified. If two of them apply, try the one that names an object you can draw.

\section*{Break one hypothesis}
Read a proof twice. The first time, cover the next line and say what would make the current line useful. If you cannot, compute a small example and a boundary case, then come back. The second time, delete one hypothesis and try to break the conclusion.

On the polynomial theorem, delete ``degree at least $1$''. The constant $41$ is prime for every $n$. The conclusion fails, so the degree hypothesis is necessary. Delete ``integer coefficients'' and the difference argument stops: $s-t$ need not divide $f(s)-f(t)$ if the coefficients are not integers. The function $f(x)=\frac12 x+2$ is not prime at $x=1$, where the value is $\frac52$, so this deletion did not break the conclusion. It broke the proof. A broken proof says you may not take that step. A broken conclusion says the missing hypothesis was necessary.

The same test catches a division by zero. From $x^2=x$ you may not cancel $x$ and announce $x=1$, because the factor you cancel is $0$ when $x=0$, and $x=0$ is a solution. If you explain the step to a friend who likes to cancel, and the friend can still cancel there, the sentence that forbids the cancellation was never written down. Write it, then cancel only where it is true.

\section*{Words on a circle}
Here is a count that uses several of those questions in turn. For an integer $n\ge2$, a word $x_1x_2\cdots x_n$ has each letter equal to $0$ or $1$. Call the word \emph{legal} when no two adjacent letters are both $1$, including the pair $(x_n,x_1)$ that closes the circle. The positions are labelled. The word $001$ is different from the word $010$.

How many legal words are there? If a word is drawn uniformly from all $2^n$ binary words of length $n$, what is the probability that it is legal?

\begin{think}
For $n=2$ the legal words are $00$, $01$ and $10$. The word $11$ is forbidden, twice over, because the two positions are adjacent in both directions. For $n=3$ the legal words are $000$, $001$, $010$ and $100$. List the legal words of length $4$. Then sort them into piles, putting two words in the same pile when a rotation takes one to the other. How many words are on your list, and how many piles?
\end{think}

There are seven legal words of length $4$, and three piles. Figure~\ref{fig:field-circle} draws them. The pile of $0000$ has one word: every rotation is the same word. The pile of $0001$ has four words. The pile of $0101$ has two, namely $0101$ and $1010$. The count of words comes first. The piles are a different question, and the natural shortcut for them is wrong.

\begin{figure}[ht]
\centering
\fig{field-circle}
\caption{The seven legal words of length $4$. A rust disc is a $1$, read clockwise from the top. Rotation sorts the seven words into piles of size $1$, $4$ and $2$. The number of piles is $3$, not $7/4$.}
\label{fig:field-circle}
\end{figure}

One language is the line count from Chapter~\ref{ch:counting}. Let $F_0=0$, $F_1=1$ and $F_{m+2}=F_{m+1}+F_m$, and write $f_m$ for the number of subsets of $\{1,\dots,m\}$ with no two consecutive elements, including the empty subset. That chapter proves $f_m=F_{m+2}$ for every $m\ge0$, with the convention $f_0=1$.

\begin{propositionx}[legal words, by the first letter]
For every integer $n\ge3$, the number of legal words of length $n$ is $F_{n+1}+F_{n-1}$.
\end{propositionx}

\begin{proof}
Split on the letter $x_1$.

Suppose $x_1=0$. The remaining letters form a word $x_2\cdots x_n$ with no two consecutive ones. The circular pairs $(x_n,x_1)$ and $(x_1,x_2)$ are automatic, because $x_1=0$. On the line of $n-1$ positions labelled $2,\dots,n$, the condition is exactly that the set of ones is a non-consecutive subset. There are $f_{n-1}=F_{n+1}$ such words.

Suppose $x_1=1$. Then $x_2=0$ and $x_n=0$, otherwise a circular pair of ones appears. The letters $x_3\cdots x_{n-1}$ are an arbitrary legal word on a line of $n-3$ positions. For $n=3$ that line is empty, and there is $f_0=1$ way to fill it, giving the single word $100$. For $n>3$ there are $f_{n-3}=F_{n-1}$ ways. In both subcases the count is $F_{n-1}$.

The two cases are disjoint and cover every legal word. The total is $F_{n+1}+F_{n-1}$.
\end{proof}

For $n=4$ the formula gives $F_5+F_3=5+2=7$, which matches the list. For $n=3$ it gives $F_4+F_2=3+1=4$. The case $n=2$ is not in the proposition: the two neighbours of position $1$ are the same position. The list at the start of the section has three legal words, and $F_3+F_1=2+1=3$ happens to agree, but that agreement is a check, not a case of the proof.

The line argument sees the circle as a line plus one awkward letter. A second language treats every position the same way. The only memory the rule needs, while you write the word from left to right, is the letter you have just written. From $0$ you may write $0$ or $1$. From $1$ you may write only $0$. Figure~\ref{fig:field-transitions} is that memory.

\begin{figure}[ht]
\centering
\fig{field-transitions}
\caption{The memory of a legal word. From $0$ the next letter may be $0$ or $1$. From $1$ it may only be $0$. A closed walk of $n$ steps, with the starting letter fixed as position $1$, is a legal circular word.}
\label{fig:field-transitions}
\end{figure}

Index the rows and columns of a $2\times2$ matrix by the letters $0,1$, in that order, and put a $1$ where a step is allowed:
\[
M=\begin{pmatrix}1&1\\1&0\end{pmatrix}.
\]
The entry $M_{ij}$ is $1$ when a letter $i$ may be followed by a letter $j$. This is the matrix of Exercise~\ref{ex:fibmat}. The new fact is what its powers count.

\begin{propositionx}[legal words, by closed walks]
For every integer $n\ge1$, the entry $(M^n)_{ij}$ equals the number of walks of $n$ steps from letter $i$ to letter $j$ along the allowed arrows. For $n\ge2$, the number of legal circular words of length $n$ is $\operatorname{tr}(M^n)$.
\end{propositionx}

\begin{proof}
For $n=1$ the entries of $M$ are the one-step walks, by the way $M$ was filled in.

Suppose the claim holds for some $r\ge1$. Matrix multiplication gives
\[
(M^{r+1})_{ij}=\sum_{k\in\{0,1\}}(M^r)_{ik}\,M_{kj}.
\]
A walk of $r+1$ steps from $i$ to $j$ has a penultimate letter $k$. The first $r$ steps form a walk from $i$ to $k$, and the last step is an allowed arrow from $k$ to $j$. Summing over the two possible letters $k$ counts those walks once each, with no overlap. Induction gives the claim for every $n\ge1$.

Now let $n\ge2$, and let $x_1\cdots x_n$ be a binary word. The word is legal if and only if each step $x_1\to x_2\to\cdots\to x_n\to x_1$ is an allowed arrow. That is a closed walk of $n$ steps whose starting letter is $x_1$. The positions stay labelled: we do not rotate the word, and position $1$ is the start. Summing $(M^n)_{ss}$ over the two possible starts $s$ counts these closed walks. That sum is the trace.
\end{proof}

Exercise~\ref{ex:fibmat} gives the entries, for every $n\ge1$,
\[
M^n=\begin{pmatrix}F_{n+1}&F_n\\F_n&F_{n-1}\end{pmatrix}.
\]
The trace is $F_{n+1}+F_{n-1}$. For $n\ge3$ this is the same count as the first-letter argument. For $n=2$ the trace is $F_3+F_1=3$, which matches the list. For $n=1$ the trace is $F_2+F_0=1$, a closed walk based at $0$ along the loop, and there is no loop at $1$. A one-letter circle was never in the problem. The trace still has a value there. Keep that value for the recurrence, and do not call it a count of words.

Define $L_0=2$, $L_1=1$ and $L_n=F_{n+1}+F_{n-1}$ for $n\ge2$. Then $L_2=3$ and $L_3=4$. For $n\ge3$,
\begin{align*}
L_{n-1}+L_{n-2}
&=(F_n+F_{n-2})+(F_{n-1}+F_{n-3})\\
&=(F_n+F_{n-1})+(F_{n-2}+F_{n-3})\\
&=F_{n+1}+F_{n-1}=L_n,
\end{align*}
where the Fibonacci recurrence was used on the pairs $(F_n,F_{n-1})$ and $(F_{n-2},F_{n-3})$, and the second pair needs $n-3\ge0$. The same relation $L_2=L_1+L_0$ holds by the numbers $3=1+2$. So $L_{n}=L_{n-1}+L_{n-2}$ for every $n\ge2$. The sequence with these initial values is the Lucas sequence. For $n\ge2$, the term $L_n$ is the number of legal circular words.

The closed form is the same comparison of initial values that gives Binet's formula after Exercise~\ref{ex:fibsums}. Set
\[
\varphi=\frac{1+\sqrt5}2,\qquad\psi=\frac{1-\sqrt5}2.
\]
Both numbers satisfy $t^2=t+1$. Indeed
\[
\varphi^2=\frac{1+2\sqrt5+5}4=\frac{3+\sqrt5}2=\varphi+1,
\]
and the calculation for $\psi$ is the same with the sign of $\sqrt5$ changed. Multiply by $\varphi^{n-2}$ for $n\ge2$: $\varphi^n=\varphi^{n-1}+\varphi^{n-2}$. The same identity holds for $\psi$. Also $\varphi^0+\psi^0=2=L_0$ and $\varphi+\psi=1=L_1$. The sequence $\varphi^n+\psi^n$ therefore satisfies the same recurrence and the same two initial values as $L_n$. By induction on $n$,
\[
L_n=\varphi^n+\psi^n
\]
for every integer $n\ge0$. The inductive step uses the recurrence, which looks two indices back, so the two initial values are the two base cases.

One numerical test is worth doing by hand. The recurrence produces
\[
L_0,\dots,L_{10}=2,1,3,4,7,11,18,29,47,76,123.
\]
Each term is the sum of the two before it. A decimal expansion of $\varphi^{10}$ begins $122.991$. The integer nearest to that decimal is $123$. Nine additions prove $L_{10}=123$. The decimal checks the additions. It is not a second proof. The bound in the next lines is why the nearest integer had to be $L_n$.

Since $\varphi\psi=(1-5)/4=-1$ and $\psi<0<\varphi$, one has $|\psi|=1/\varphi$. Also $\varphi>1$, because $(1+\sqrt5)/2>1$ is $\sqrt5>1$. Thus $|\psi|<1$, and $|\psi|<\varphi$. For $n\ge2$,
\[
|\psi|^n\le|\psi|^2=\frac{3-\sqrt5}2.
\]
The comparison $|\psi|^2<\frac12$ is $3-\sqrt5<1$, that is $\sqrt5>2$, that is $5>4$. Thus $|L_n-\varphi^n|=|\psi|^n<\frac12$ for every $n\ge2$, and $L_n$ is the integer closest to $\varphi^n$.

Under the uniform probability on the $2^n$ binary words of length $n\ge2$, the probability of a legal word is
\[
p_n=\frac{L_n}{2^n}=\Bigl(\frac{\varphi}2\Bigr)^n+\Bigl(\frac{\psi}2\Bigr)^n.
\]
If $0<q<1$, write $1/q=1+h$ with $h>0$. Bernoulli's inequality, Proposition~\ref{prop:bernoulli}, gives $(1+h)^n\ge1+nh$. Given $\varepsilon>0$, choose an integer $N$ with $1+Nh\ge1/\varepsilon$. Then $q^n\le\varepsilon$ for every $n\ge N$, so $q^n\to0$.

The two geometric terms in $p_n$ are not the same size. From $|\psi|=1/\varphi$ one has $|\psi/\varphi|=1/\varphi^2$. Since $\varphi>1$, this ratio lies strictly between $0$ and $1$, so $|\psi/\varphi|^n\to0$ by the decay just proved. The signed powers $(\psi/\varphi)^n$ have the same absolute value, and therefore tend to $0$ as well. Hence
\[
\frac{L_n}{\varphi^n}=1+\Bigl(\frac{\psi}{\varphi}\Bigr)^n\to1.
\]
Chapter~\ref{ch:sequences} writes $u_n\sim v_n$ when $u_n/v_n\to1$. In that sense $L_n\sim\varphi^n$, and $p_n\sim(\varphi/2)^n$. The equivalence compares the exact count with its leading term. It does not by itself show that $p_n$ tends to $0$. For the limit, use $\varphi<2$: the inequality $(1+\sqrt5)/2<2$ is $\sqrt5<3$, that is $5<9$. Combined with $|\psi|<\varphi$,
\[
0<p_n<2\Bigl(\frac{\varphi}2\Bigr)^n.
\]
The base $\varphi/2$ lies strictly between $0$ and $1$, so the right side tends to $0$. The squeeze in Chapter~\ref{ch:sequences} gives $p_n\to0$.

For a long word, drawn uniformly, the probability of being legal tends to $0$. The exact count $L_n$ still grows. The set of all binary words grows faster, because $\varphi<2$. A table of the first twenty probabilities would be a list. Both the formula for $p_n$ and the limit are laws.

\begin{beauty}
The same integer $L_n$ is a trace, a sum of two Fibonacci numbers, and the integer closest to $\varphi^n$. The trace counts closed walks. The Fibonacci sum counts non-consecutive subsets of a line, after one letter has been peeled off and handled separately. The power $\varphi^n$ is what remains once $\psi^n$ is smaller than $\frac12$. The probability is that integer divided by $2^n$. It tends to $0$ because $\varphi<2$, while $L_n$ itself keeps growing.
\end{beauty}

\subsection*{Two shortcuts that fail}
Dividing $L_n$ by $n$ does not count the piles. For $n=4$, the quotient $7/4$ is not an integer. A pile of words cannot have size $7/4$. The division would be right if every pile contained $n$ distinct words, one for each rotation. Figure~\ref{fig:field-circle} shows that the piles have sizes $1$, $4$ and $2$. The word $0000$ is fixed by every rotation. The word $0101$ comes back after two turns. Only $0001$ has four distinct rotations. The hidden hypothesis was that every pile has the same size. It is false, so the quotient is not the number of piles. There are three piles. Finding a formula for the number of piles, for a general $n$, is a different problem. It starts from the hypothesis that just failed, and this volume stops there.

A second shortcut multiplies probabilities that are not independent. Any fixed adjacent pair, including the closing pair, is one of $00,01,10,11$ with equal probability, so the chance that it is not $11$ is $3/4$. The product $(3/4)^n$ is the probability you would get if those $n$ events were independent. They are not: consecutive pairs share a letter, and the closing pair shares its letters with the first pair. The numbers disagree immediately. For $n=2$, the product is $9/16$, while $p_2=3/4$. For $n=3$, the product is $27/64$, while $p_3=4/8=1/2$. Chapter~\ref{ch:probability} allows the product only under an independence hypothesis you have checked. Here the exact count is the justified step, and the product is not.

Both failures were found by the same question. What does the shortcut assume, and is there a single $n$ at which that assumption is false? One counterexample ends it.

\section*{A decimal is a check}
A calculation can catch a false identity, and it can catch an arithmetic slip in a true one. The expansion of $\varphi^{10}$ beginning $122.991$ is a check that the tenth Lucas number really is $123$. If the recurrence had produced $124$, the decimal would have said to add the list again. The same calculation cannot promote a finite pattern to a universal sentence. A machine that tested $f(n)$ for $0\le n\le39$ and reported forty primes would be repeating Euler's check. It would not have proved a theorem about every $n$, and it would have missed the factorisation $f(40)=41^2$ unless it was asked for $n=40$.

Three kinds of evidence stay different, and they should keep their names. A symbolic identity, such as $f(40)=41^2$ or $L_n=\varphi^n+\psi^n$, is an equality of exact expressions. A truncation bound, such as $|L_n-\varphi^n|<\frac12$, or the Taylor remainder in the hat check, says how far an approximation is allowed to be. A decimal, or a picture of the first few words, is a witness for the values it actually computed. Use all three. Do not treat one of them as if it were another.

\section*{From a finite list of primes to a law}
The matrix counted words. It did not see that $41$ divides $f(40)$. The theorem on prime values of a polynomial used that kind of divisibility and stopped. It did not say whether the primes themselves go on forever, or how large a gap between them can be. Those sentences are laws about the integers. They need a theory of divisibility. Another identity of the shape $n^2+n+41$ will not prove them.

The next book is Volume~2, \emph{The Integers: Primes, Patterns and Proofs}. It takes the division algorithm, B\'ezout's identity and the unique factorisation proved in Section~\ref{sec:divisibility} of Chapter~\ref{ch:sums}, ``Divisibility from well ordering,'' and it makes them the beginning of the subject. The object is a prime.

Here is a picture of the first question, small enough to compute. The product of the first six primes is
\[
2\cdot3\cdot5\cdot7\cdot11\cdot13=30030,
\]
so the next integer is $30031$. Check the factorisation $59\cdot509$: one has $59\cdot500=29500$ and $59\cdot9=531$, and $29500+531=30031$. The factor $59$ is prime. It is not divisible by $2$, $3$, $5$ or $7$, and the next prime is $11$, whose square $121$ is already larger than $59$. And $59$ is not one of the six primes that were multiplied. So $30031$ is composite, and it still has a prime factor off the list. This single number is not a proof for every finite list. It is the object the proof uses. Volume~2 writes the argument for an arbitrary finite list of primes. The product plus one need not itself be prime. The number $30031$ was not.

The same book then proves a sharper existence statement. Bertrand's postulate, as proved there: for every integer $n\ge1$ there is a prime $p$ with $n<p\le2n$. The argument compares sizes inside the binomial coefficient $\binom{2n}{n}$. You already know that coefficient as a count. The new work is to see what its size forces on the primes.

You still write the sentence, try a number you can finish, and name the hypothesis a shortcut would need. The object changes. A legal word was a closed walk. A prime, as defined in Chapter~\ref{ch:sums}, is an integer greater than $1$ whose only positive divisors are $1$ and itself. Volume~2 proves that such integers do not run out, and it is where the patterns of the integers start to carry proofs of their own.
